\section{Step 68 Publish a complete block result}\label{h263-block-s08}
After every dependency and pixel operation succeeds, copy the selected levels,
reconstructed coefficients, inverse values and final pixels into the caller's
448-byte record. Return the saved ordinary count. A child can fail after writing
its own private output; none of that work authorizes an early public store.

Stop the ordinary caller at \texttt{block\_caller\_return}, before its first
instruction. The complete record and return registers are ready, but named
status/count fields still contain their old values. The default prints first
pixel 131 and exits zero after the caller resumes.

\paragraph{Observe} The default successful count is one. INTRA input 128 instead
produces DC code 255, reconstructed DC 1024, pixels 128 and ordinary count zero.
\Needspace{8\baselineskip}
\paragraph{Check} What survives a failed inverse child?
\begin{itemize}
  \item[A.] A partially replaced caller record.
  \item[B.] Every old caller-result byte, with parent failure status two.
  \item[C.] A successful status because forward transform already completed.
\end{itemize}
